% =====================================================================
% CHAPTER 3: DATA AND SIMULATION DESIGN                     budget: 6 pages
% ROLE: Describe the real data and the simulator well enough that a
%       reader could rebuild them. No results here.
% =====================================================================
\chapter{Data and Simulation Design}
\label{ch:design}

\section{Enron data and hierarchy labels}
% budget: 1.5 pages

% P1 KEY: How the raw corpus was turned into usable messages and threads.
% SUPPORT:
%   - Rebuild the message column, parse headers, find threads (Exposé p.8).
%   - TODO: fill in from PROGRESS.md / PROGRESS_llms.md (preprocessing
%     steps, duplicates removed, final counts).

% P2 KEY: Rank comes from job titles. Direction comes from comparing the
%         ranks of sender and recipient.
% SUPPORT:
%   - Job titles of 156 employees, 6 ranks; departments Trading and Legal.
%   - Direction: down / to a peer / up.
%   - Limitation (formal rank only): see Ch. 2 P1 and Ch. 8.
%   - SRC: RO 2.0.

% P3 KEY: Which real-email samples are used, and for what.
% SUPPORT:
%   - Q1 benchmark: 2,202 emails from 107 senders.
%   - Design A: 183 real threads (from 200 sampled).
%   - Direction judge (Layer 1): 1,001 real emails.
%   - Tone and politeness: 750 per source, 250 per direction.
%   - TAB: sample, size, used in which section.

\section{Models}
% budget: 1 page

% P4 KEY: Four models of different size and family, all free.
% SUPPORT:
%   - TAB: model, family, where it runs, role (writer / judge).
%   - Llama 3.2 3B local (Ollama); gpt-oss-20b, gpt-oss-120b, qwen3.8-27b
%     on NVIDIA and Groq free tiers.
%   - Why: no paid API; the laptop cannot run models above about 8B.
%   - SRC: RO 0 (Models).

% P5 KEY: Free tiers set practical limits, and a cache makes every result
%         reproducible.
% SUPPORT:
%   - Groq: about 1,000 requests and 200,000 tokens per day.
%   - Every reply is cached locally; a result never depends on the service.
%   - Total cost $0.
%   - Prompts with Enron text go to third-party servers (link to ethics, Ch. 8).

\section{Prompt structure}
% budget: 1.5 pages
% FIG: system diagram (update the Exposé diagram, p.9, to what was built).

% P6 KEY: Each persona is a role, built from real writing statistics, and
%         never a named person.
% SUPPORT:
%   - 10 personas: 5 ranks x 2 departments (Trading, Legal).
%   - Each summarizes how real employees at that rank wrote: length,
%     recipients, share of orders, hedges, deference, questions.
%   - Each built from at least 3 real people; rank 6 left out (too few).
%   - No names: the model may have memorized Enron (Ch. 2 P9).
%   - SRC: RO 1.1.

% P7 KEY: Each persona has a small memory, following Generative Agents.
% SUPPORT:
%   - 15 invented recent events (observations) and a few reflections.
%   - The most relevant ones for persona and direction go in the prompt.
%   - Park et al. 2023.

% P8 KEY: The reply comes as structured output. Its decision field is the
%         "decision attitude" of Q1.
% SUPPORT:
%   - 5 fields: subject, body, decision (accept / decline / defer /
%     escalate / none), confidence (low / medium / high), reason.

% P9 KEY: A hand reading found that small models hand requests back to the
%         sender, so the prompt says the persona must act.
% SUPPORT:
%   - About a quarter of 100 Llama replies mirrored the request.
%   - The added paragraph describes the behavior, not words, so it does not
%     target the mirroring measure.
%   - On Llama it changed little (borrowed words 0.579 -> 0.565, p=.61).
%   - Full prompt in the appendix.

\section{Three generation designs}
% budget: 1.5 pages
% TAB: the three designs side by side (input, personas, size, used for).

% P10 KEY: Design A answers real Enron emails, so each generated reply has a
%          real reply to compare with (Q2).
% SUPPORT:
%   - First message of a real thread; the persona matches the real replier's
%     rank and department; direction from the real ranks.
%   - 183 threads. Same paired logic as HEART.
%   - SRC: RO 1.2.

% P11 KEY: Design B is a full scenario grid that varies direction under
%          controlled conditions (Q1, Q3).
% SUPPORT:
%   - 6 task types x 3 directions x 2 stakes x 4 incoming tones = 144 scenarios.
%   - x 10 personas = 1,440 replies per draw. Llama: 3 draws (4,320).
%   - Incoming tone is a property of the message, not an instruction.
%   - Judge-swap set: 60 scenarios (2 task types x 3 directions, neutral tone).
%   - SRC: RO 1.3.

% P12 KEY: Design C uses 40 personas, so a rank effect cannot come from
%          one persona's prompt.
% SUPPORT:
%   - Problem: 10 personas = 10 clusters; clustered errors too small
%     (Cameron & Miller 2015); Llama's p moved from .248 to .018 between draws.
%   - 4 variants per role: tenure (2 or 12 years) x origin (internal / external).
%   - Each writes to a named receiver one rank up, same rank, one rank down.
%   - Pilot: 832 replies per model; qwen stopped (Groq limit).
%   - SRC: RO 1.4.

% P13 KEY: Generation worked for all models at no cost, but models write
%          shorter replies than people.
% SUPPORT:
%   - TAB: generation summary (RO 1.5): valid replies, mean length.
%   - Design A mean words: Llama 22.7, 20b 32.8, 120b 28.6, qwen 48.3, real 65.0.
%   - Length matters for several measures (Ch. 4, length matching).

\section{Changes from the original plan}
% budget: 0.5 pages

% P14 KEY: Some Exposé parts were dropped or changed, each for a reason.
% SUPPORT:
%   - Power score per message -> did not track rank; replaced by
%     sentence-level orders and judges.
%   - Scenario builder from threads -> Design A uses threads directly.
%   - Judge from another model family -> not possible on the free path;
%     replaced by the judge-swap design (Ch. 7).
%   - Human coders -> TODO: done or not? Decide before writing Ch. 7.
